% German abstract (Zusammenfassung), a translation of the English abstract of thesis_v3_updated.tex.
% NOT included in the thesis. Only if the exam office requires a German abstract: insert the block
% below right after \end{abstract} in thesis_v3_updated.tex (before \setcounter{page}{2}), then rebuild.
% It needs no extra package (tested 25 Nov 2026: compiles, no overfull lines).
% Every figure is the same as in the English abstract; if that abstract changes, change this too.

{\renewcommand{\abstractname}{Zusammenfassung}
\begin{abstract}
\thispagestyle{fancy}
Ein Katalogisierungssystem für ein historisches Archiv ist nur dann nützlich, wenn es angeben kann, welchen seiner Antworten ohne Prüfung durch einen Menschen vertraut werden kann. Diese Arbeit untersucht diese Frage an 12.110 Illustrationen aus historischen Büchern, für zwei Aufgaben, deren Klassenlisten zu Beginn festgelegt wurden: welche von fünf Szenenklassen (Unterricht, Familie, Spiel, Landschaft, Zeichnung) ein Bild zeigt, falls überhaupt, und welche von acht Objektklassen es enthält. Drei Modelle mit unterschiedlicher Architektur und unterschiedlichen Trainingsdaten dienen als Agenten: ein Objektdetektor, ein Zero-Shot-Bild-Text-Abgleich und ein Vision-Language-Modell. Das System beruhte auf der Hypothese, dass ihre Uneinigkeit auf Zweifel hinweist; seine Entscheidungen werden mit Conformal Risk Control gesteuert, das den Anteil der automatisch freigegebenen und falschen Antworten begrenzt.

Die Hypothese trifft nicht zu. In einem vorab festgelegten Test ordnet die Uneinigkeit die Fehler des Systems mit einer AUROC von $0{,}378$, die eigene Konfidenz des Systems mit $0{,}825$; innerhalb jeder Aufgabe gemessen erreicht die Uneinigkeit $0{,}78$ für Objektlabels und $0{,}82$ für das Vorliegen einer Szene, liegt aber noch immer nicht über der Konfidenz ($0{,}84$ und $0{,}83$): Die meisten Fehler betreffen Bilder, bei denen alle Modelle gemeinsam falschliegen. Ein Detektor, der an den Pseudo-Labels seines Lehrermodells gemessen wird, erreicht $0{,}768$, an menschlichen Boxen aber $0{,}146$, weil der Lehrer die meisten Klassen nie gelabelt hat; die eigene Auswertung der Arbeit bewertet dieselben Detektionen mit $0{,}123$, und ein besserer Lehrer hebt diesen Wert auf $0{,}210$.

Die Arbeit schlägt deshalb ein korrigiertes Design vor, in dem jedes Modell berichtet, was es misst, und die Konfidenz entscheidet, was ein Mensch prüft. Für die acht Objektklassen ordnen die Modelle das Vorhandensein mit einer mittleren AUROC von $0{,}903$, und ein Versprechen über übersehene Objekte gibt jedem Bild eine Liste von $2{,}8$ Klassen, die etwa eines von zehn Objekten auslässt ($10{,}8\,\%$ aller Objekte). Für die fünf Szenenklassen erstellen sie eine Vorauswahl, können aber nicht entscheiden: Unter einem Risikoversprechen von $5\,\%$ wird ,,keine Szene`` zuverlässig freigegeben, ,,Zeichnung`` mit etwa einer falschen von sechs Antworten, und die spezifischen Klassen brauchen einen Menschen; auf zufälligen Archivbildern, auf denen Spielszenen häufiger sind als Zeichnungen, fällt die balancierte Genauigkeit von $0{,}544$ auf $0{,}351$. Ob ein Bild überhaupt eine Szene zeigt, wird bei einem Versprechen von $5\,\%$ für $76{,}4\,\%$ von 1.250 gelabelten Bildern freigegeben, gegenüber $63{,}0\,\%$ für das ursprünglich eingesetzte System. Auf 400 zufälligen Archivbildern scheitert dieses Versprechen in einem vorab festgelegten Test ($16{,}2\,\%$ freigegeben und falsch), weil die Labels zwei Definitionen einer Szene vermischten und die gelabelten Bilder keine Zufallsstichprobe sind. Mit einer Definition und einer Kalibrierung auf zufälligen Archivbildern liegt das Risiko knapp über dem versprochenen Niveau ($5{,}4$--$5{,}9\,\%$ auf 294 neuen Bildern). Ein einzelnes Vision-Language-Modell ist etwa so gut wie die Kombination; für die Szenenklassen muss die Kombination gelernt statt festgelegt werden. Der Beitrag der Arbeit ist eine geprüfte Antwort darauf, wann einer automatischen Antwort vertraut werden kann: nach Konfidenz steuern, nicht nach Übereinstimmung zwischen Modellen; auf einer Zufallsstichprobe der Sammlung selbst kalibrieren; und versprechen, was jede Aufgabe braucht: eine Schranke für übersehene Objekte oder eine kurze Liste von Szenenklassen.
\end{abstract}}
